\documentclass[UTF8,fontset=none,11pt,a4paper]{ctexart}
\usepackage{ipara-handbook}
\handbooksetup{DS-06 Union-Find 与集合划分}{408 · 数据结构 DS-06}{Partition · Representative · Find · Union}
\begin{document}
\handbooktitle{Union-Find 与集合划分}{动态合并后，怎样快速回答“是否同一集合”}{408 · 数据结构 Topic DS-06}

\begin{corebox}{母问题：动态连通不是一般树，而是代表元问题}
并查集的生成链是
\[
\text{Partition}\to\text{Representative}\to\text{Find/Union}\to\text{Path Compression + Rank}\to\text{Dynamic Connectivity}.
\]
每个集合用一棵父指针树表示，但树形只是实现手段；真正的抽象是分区。\code{find(x)} 返回集合代表元，\code{union(a,b)} 合并两个代表元不同的集合。任何操作后，父指针必须最终通向唯一根，根的数量必须等于集合数量。
\end{corebox}
\begin{methodbox}{本册定位与 Owner 边界}
\begin{itemize}
\item \kw{Owns：}MakeSet、Find、Union、代表元、按秩/大小合并、路径压缩与动态连通语义。
\item \kw{Uses：}数组下标和成本向量；向 DS-B03/DS08 输出 `connected/unite` 接口。
\item \kw{Stops：}一般树遍历、Kruskal 的边权排序与生成树正确性由其他 Owner 负责。
\end{itemize}
\end{methodbox}
\tableofcontents\newpage

\section{分区不变量}
初始时每个元素自成集合，$parent[i]=i$。父指针反复追踪必到根；根是自己的父节点。若两个元素根相同，则属于同一集合；若根不同，合并时把一个根挂到另一个根下，并将集合数减一。重复合并不改变分区。

\section{为什么需要两种局部修复}
若只把一个根挂到另一个根下，连续合并可能生成长链，使 Find 退化到 $O(n)$。按秩合并总把较矮树挂到较高树，限制树高；路径压缩在 Find 返回时让沿途节点直接指向根，摊还复杂度达到近似常数 $O(\alpha(n))$。

\section{代码：代表元、路径压缩与按秩合并}
完整实现位于 \codepath{code/ds06_union_find.hpp}，测试位于 \codepath{code/ds06_union_find_test.cpp}；正文直接引用真实源码。
\lstinputlisting[language=C++,firstline=11,lastline=58,firstnumber=11,numbers=left,caption={Find 的路径压缩与 Union 的按秩合并}]{30_408/10_数据结构/06_UnionFind与集合划分/code/ds06_union_find.hpp}
`find` 的递归返回值就是代表元；回溯赋值是路径压缩。`unite` 先分别 Find，再比较秩，最后只减少一次 `components`。根相同的分支必须提前返回，否则重复合并会错误减少集合数。
\lstinputlisting[language=C++,firstline=60,lastline=82,firstnumber=60,numbers=left,caption={分区不变量和越界边界}]{30_408/10_数据结构/06_UnionFind与集合划分/code/ds06_union_find.hpp}

\subsection{测试：重复合并、压缩与空全集}
\lstinputlisting[language=C++,firstline=7,lastline=52,firstnumber=7,numbers=left,caption={代表元一致性和集合数变化测试}]{30_408/10_数据结构/06_UnionFind与集合划分/code/ds06_union_find_test.cpp}
\begin{lstlisting}[language=bash]
clang++ -std=c++17 -Wall -Wextra -Werror -pedantic \\
  code/ds06_union_find_test.cpp -o /tmp/ds06_test
/tmp/ds06_test
\end{lstlisting}

\section{复杂度与边界}
\begin{center}\small
\begin{tabularx}{\linewidth}{L{2.4cm}C{2cm}YY}\toprule
操作 & 摊还复杂度 & 保证 & 边界\\\midrule
Find/Union & $O(\alpha(n))$ & 路径压缩 + 按秩/大小 & 空集合、越界\\
MakeSet & $O(n)$ & 每项自为根 & $n=0$ 合法\\
空间 & $O(n)$ & parent/rank 两数组 & 不保存集合内部顺序\\
\bottomrule\end{tabularx}\end{center}
\begin{warnbox}{不要把“父指针树”当作业务树}
并查集不支持按层遍历、输出路径或保持关键字有序；它只回答代表元和同集合关系。把 Union-Find 的父树画成一般树来做遍历，会把实现结构误当成目标对象。
\end{warnbox}

\section{做题控制协议与压缩复原}
先写当前分区和每个根；每次合并只比较两个代表元；重复合并检查是否同根；若问复杂度，指出路径压缩和秩策略共同作用。忘掉代码后复原：\kw{集合如何表示？谁是代表元？Find 是否压缩？Union 挂哪棵树？集合数何时减少？}
\begin{corebox}{跨册交接}
DS06 向 DS-B03 和 DS08 只提供 `connected/unite`。Kruskal 的排序、选边和无环性证明由 DS08 Own；本册只保证每次加入边前后动态分区不变量。
\end{corebox}
旧总册关于代表元、路径压缩、按秩合并和 Kruskal 接口的内容已吸收；带权并查集、可回滚并查集和持久化版本作为 Extension。
\end{document}
